\MathBookSourcePage{156}
\subsection{一般椭圆问题的变分逼近}
\label{sec:ch05-general-elliptic-variational-approximation}
\setcounter{thmcount}{0}
\setcounter{equation}{0}

上一节看到, 某些适定的椭圆问题所对应的变分问题不具有强制性, 但加上适当大的常数总能使它具有强制性. 在这种情形下, 遵循 Schatz 的思想(Schatz 1974), 仍可给出变分问题有限元逼近的估计, 只是论证稍复杂. 首先假设变分形式 $a(\cdot,\cdot)$ 满足下列性质. 假设 $a(\cdot,\cdot)$ 在 $H^1(\Omega)$ 上连续:
\MBSourceItem{1}
\begin{equation}
\MathBookFormulaID{formula-ch05-general-form-continuity}
\label{eq:ch05-general-form-continuity}
|a(u,v)|\leq C_1\|u\|_{H^1(\Omega)}\|v\|_{H^1(\Omega)}
\quad\forall u,v\in H^1(\Omega),
\end{equation}
并假设某个附加常数 $K\in\mathbb R$ 使它具有强制性:
\MBSourceItem{2}
\begin{equation}
\MathBookFormulaID{formula-ch05-shifted-coercivity}
\label{eq:ch05-shifted-coercivity}
a(v,v)+K(v,v)\geq\alpha\|v\|_{H^1(\Omega)}^2
\quad\forall v\in H^1(\Omega).
\end{equation}
假设存在某个 $V\subset H^1(\Omega)$, 使变分问题
\[
\MathBookFormulaID{formula-ch05-general-continuous-primal-problem}
a(u,v)=(f,v)\quad\forall v\in V
\]
以及伴随变分问题
\[
\MathBookFormulaID{formula-ch05-general-continuous-adjoint-problem}
a(v,u)=(f,v)\quad\forall v\in V
\]
均有唯一解 $u$, 且两种情形都对所有 $f\in L^2(\Omega)$ 满足正则性估计
\MBSourceItem{3}
\begin{equation}
\MathBookFormulaID{formula-ch05-general-h2-regularity}
\label{eq:ch05-general-h2-regularity}
|u|_{H^2(\Omega)}\leq C_R\|f\|_{L^2(\Omega)}.
\end{equation}

令 $V_h$ 是第~\ref{sec:ch05-poisson-variational-approximation} 节所述的 $V$ 的有限元子空间, 并由
\MBSourceItem{4}
\begin{equation}
\MathBookFormulaID{formula-ch05-general-discrete-problem}
\label{eq:ch05-general-discrete-problem}
a(u_h,v)=(f,v)\quad\forall v\in V_h
\end{equation}
定义 $u_h\in V_h$. 表面看来, 式~\eqref{eq:ch05-general-discrete-problem} 未必有唯一解, 因为 $a(\cdot,\cdot)$ 可能不具有强制性. 实际上无法保证它对所有子空间 $V_h$ 都有唯一解. 但假设
\MBSourceItem{5}
\begin{equation}
\MathBookFormulaID{formula-ch05-general-approximation-property}
\label{eq:ch05-general-approximation-property}
\inf_{v\in V_h}\|u-v\|_{H^1(\Omega)}
\leq C_Ah|u|_{H^2(\Omega)}
\quad\forall u\in H^2(\Omega).
\end{equation}
下面的结果应与 Céa 定理~\ref{thm:ch02-cea} 比较.

\MBSourceItem{6}
\begin{Theorem}
\label{thm:ch05-schatz-error-estimate}
在条件~\eqref{eq:ch05-general-form-continuity}、\eqref{eq:ch05-shifted-coercivity}、\eqref{eq:ch05-general-h2-regularity} 和~\eqref{eq:ch05-general-approximation-property} 下, 存在常数 $h_0,C$, 使对所有 $h\leq h_0$, 式~\eqref{eq:ch05-general-discrete-problem} 有唯一解, 且
\[
\MathBookFormulaID{formula-ch05-schatz-quasi-optimality}
\|u-u_h\|_{H^1(\Omega)}
\leq C\inf_{v\in V_h}\|u-v\|_{H^1(\Omega)},
\]
其中可取 $C=2C_1/\alpha$, 并且
\[
\MathBookFormulaID{formula-ch05-schatz-l2-error}
\|u-u_h\|_{L^2(\Omega)}
\leq C_1C_AC_Rh\|u-u_h\|_{H^1(\Omega)}.
\]
特别地, 可取
\[
\MathBookFormulaID{formula-ch05-schatz-mesh-threshold}
h_0=(\alpha/2K)^{1/2}/(C_1C_AC_R).
\]
\end{Theorem}
\begin{Proof}
先对任意可能存在的式~\eqref{eq:ch05-general-discrete-problem} 的解导出估计. 无论如何总有
\MathBookSourcePage{157}
\MBSourceItem{7}
\begin{equation}
\MathBookFormulaID{formula-ch05-galerkin-orthogonality}
\label{eq:ch05-galerkin-orthogonality}
a(u-u_h,v)=0\quad\forall v\in V_h.
\end{equation}
由式~\eqref{eq:ch05-shifted-coercivity}、\eqref{eq:ch05-galerkin-orthogonality} 和~\eqref{eq:ch05-general-form-continuity}, 对任意 $v\in V_h$,
\[
\MathBookFormulaID{formula-ch05-schatz-energy-bound}
\begin{aligned}
\alpha\|u-u_h\|_{H^1(\Omega)}^2
&\leq a(u-u_h,u-u_h)+K(u-u_h,u-u_h)\\
&=a(u-u_h,u-v)+K\|u-u_h\|_{L^2(\Omega)}^2\\
&\leq C_1\|u-u_h\|_{H^1(\Omega)}\|u-v\|_{H^1(\Omega)}
+K\|u-u_h\|_{L^2(\Omega)}^2.
\end{aligned}
\]
使用标准对偶性方法估计 $\|u-u_h\|_{L^2(\Omega)}$. 令 $w$ 为条件~\eqref{eq:ch05-general-h2-regularity} 所保证的伴随问题的解:
\[
\MathBookFormulaID{formula-ch05-schatz-adjoint-error-problem}
a(v,w)=(u-u_h,v)\quad\forall v\in V.
\]
则对任意 $w_h\in V_h$,
\[
\MathBookFormulaID{formula-ch05-schatz-duality-chain}
\begin{aligned}
(u-u_h,u-u_h)
&=a(u-u_h,w)\\
&=a(u-u_h,w-w_h)\qquad\text{(由式~\eqref{eq:ch05-galerkin-orthogonality})}\\
&\leq C_1\|u-u_h\|_{H^1(\Omega)}\|w-w_h\|_{H^1(\Omega)}.
\end{aligned}
\]
适当选择 $w_h$, 得
\[
\MathBookFormulaID{formula-ch05-schatz-duality-regularity}
\begin{aligned}
(u-u_h,u-u_h)
&\leq C_1C_Ah\|u-u_h\|_{H^1(\Omega)}|w|_{H^2(\Omega)}\\
&\leq C_1C_AC_Rh\|u-u_h\|_{H^1(\Omega)}\|u-u_h\|_{L^2(\Omega)}.
\end{aligned}
\]
所以
\MBSourceItem{8}
\begin{equation}
\MathBookFormulaID{formula-ch05-schatz-dual-error-bound}
\label{eq:ch05-schatz-dual-error-bound}
\|u-u_h\|_{L^2(\Omega)}\leq Ch\|u-u_h\|_{H^1(\Omega)},
\end{equation}
其中 $C=C_1C_AC_R$. 代回上式, 得
\[
\MathBookFormulaID{formula-ch05-schatz-absorption-bound}
\alpha\|u-u_h\|_{H^1(\Omega)}^2
\leq C_1\|u-u_h\|_{H^1(\Omega)}\|u-v\|_{H^1(\Omega)}
+KC^2h^2\|u-u_h\|_{H^1(\Omega)}^2.
\]

\MathBookSourcePage{158}
因此当 $h\leq h_0$, 其中 $h_0=(\alpha/2K)^{1/2}/(C_1C_AC_R)$, 有
\MBSourceItem{9}
\begin{equation}
\MathBookFormulaID{formula-ch05-schatz-final-energy-estimate}
\label{eq:ch05-schatz-final-energy-estimate}
\alpha\|u-u_h\|_{H^1(\Omega)}
\leq2C_1\|u-v\|_{H^1(\Omega)}
\quad\forall v\in V_h.
\end{equation}

以上一直在假设解 $u_h$ 存在. 现在考虑其存在性与唯一性. 因为式~\eqref{eq:ch05-general-discrete-problem} 是未知量个数与方程个数相同的有限维系统, 存在性与唯一性等价. 不唯一将蕴含当 $f\equiv0$ 时存在非平凡解 $u_h$. 此时由式~\eqref{eq:ch05-general-h2-regularity} 有 $u\equiv0$, 但当 $h$ 足够小时, 式~\eqref{eq:ch05-schatz-final-energy-estimate} 又蕴含 $u_h\equiv0$. 因而 $h$ 足够小时, 式~\eqref{eq:ch05-general-discrete-problem} 有唯一解. 最后, 当 $h\leq h_0$ 时, 该唯一解满足式~\eqref{eq:ch05-schatz-final-energy-estimate} 和~\eqref{eq:ch05-schatz-dual-error-bound}, 证明完成.
\end{Proof}
